%-------------------------------------------------------------------------------
%                            BAB II
%                       TINJAUAN KEPUSTAKAAN
%-------------------------------------------------------------------------------

\chapter{TINJAUAN KEPUSTAKAAN}

% ======================================================================
% 2.1 Tanah Vulkanik
% ======================================================================
\section{Karakteristik Geokimia Tanah Vulkanik}

\subsection{Proses Pembentukan dan Karakteristik Tanah Vulkanik}
% TODO: Isi dengan referensi primer tentang proses pembentukan tanah vulkanik
%       dari material piroklastik/lava, proses pelapukan, alterasi mineral,
%       serta sifat fisikokimia umum tanah hasil erupsi.

\subsection{Mineralogi Tanah Vulkanik}
% TODO: Isi dengan pembahasan mineral utama penyusun tanah vulkanik seperti
%       feldspar, olivin, piroksen, kuarsa, dan oksida besi, beserta kaitannya
%       dengan kandungan unsur mayor.

\subsection{Unsur Mayor pada Tanah Vulkanik}
% TODO: Isi dengan pembahasan unsur mayor (Si, Al, Fe, Ca, Mg, Na, K) yang
%       dominan pada tanah vulkanik, rentang konsentrasi tipikal, dan mineral pembawanya.

\subsection{Variabilitas Geokimia dan Faktor Lingkungan}
% TODO: Isi dengan pembahasan variasi komposisi unsur akibat arah deposisi abu,
%       kedalaman tanah, curah hujan, topografi, angin, dan proses pencucian unsur.

\subsection{Karakteristik Spektral Material Vulkanik}
% TODO: Isi dengan pembahasan karakteristik spektrum emisi LIBS pada material
%       vulkanik, kemungkinan overlap garis emisi antar unsur, dan kompleksitas matriks spektral.

\subsection{Peran Unsur terhadap Kesuburan dan Vulkanologi}
% TODO: Isi dengan pembahasan peran unsur mayor terhadap kesuburan tanah,
%       evaluasi lahan pertanian, serta relevansinya dalam pemantauan aktivitas
%       vulkanik dan mitigasi bencana.


% ======================================================================
% 2.2 Laser-Induced Breakdown Spectroscopy (LIBS)
% ======================================================================
\section{Laser-Induced Breakdown Spectroscopy (LIBS)}

\subsection{Prinsip Dasar LIBS}

\textit{Laser-Induced Breakdown Spectroscopy} (LIBS) adalah teknik spektroskopi emisi atomik yang memanfaatkan interaksi destruktif antara pulsa laser berintensitas tinggi dengan permukaan material untuk menentukan komposisi unsur penyusunnya. Saat pulsa laser, yang umumnya berdurasi orde nanosekon, difokuskan pada area target, energi optik diserap secara masif oleh material melalui proses pemanasan, pelelehan, dan penguapan seketika. Proses ablasi ini menghasilkan awan uap berkerapatan tinggi yang secara cepat terus menyerap sisa energi laser yang datang melalui mekanisme penyerapan seperti \textit{Inverse Bremsstrahlung} dan efek multiphoton. Kondisi pemerangkapan energi ini memicu pemecahan dielektrik (\textit{dielectric breakdown}) dari gas di sekitarnya dan membentuk plasma mikro bersuhu amat tinggi yang kaya akan molekul, atom eksitasi, ion, serta elektron bebas \citep{legnaioli_laser-induced_2025}.

\subsection{Pembentukan Plasma LIBS}

Sembari plasma tersebut berekspansi secara eksponensial dalam ruang bebas setelah berhentinya pemaparan laser, plasma mulai mendingin yang secara bertahap mengurangi pendaran emisi radiasi hamburan (\textit{continuum background}). Di saat yang sama, atom dan ion di dalam plasma melepas surplus energinya kembali ke tingkat orbit dasar melalui pancaran spektrum optik sekunder. Kemampuan LIBS untuk menganalisis material secara \textit{in situ} dan minim \textit{sample preparation} menjadikan LIBS sangat unggul untuk observasi geokimia benda padat kompleks, khususnya sampel geologi bervariasi tinggi layaknya tanah vulkanik tempat fenomena efek matriks banyak terjadi \citep{sawyers_database_2025, hao_machine_2024}.

\subsection{Emisi Atomik dan Identifikasi Unsur}
% TODO: Isi dengan pembahasan tentang mekanisme emisi atomik diskrit,
%       identifikasi unsur berdasarkan panjang gelombang karakteristik,
%       dan peran basis data NIST ASD dalam identifikasi garis emisi.

\subsection{Local Thermodynamic Equilibrium (LTE)}

Analisis kuantitatif dari emisi plasma LIBS yang mengandalkan pemodelan fisika analitik umumnya bersandar pada validitas asumsi Keseimbangan Termodinamika Lokal atau \textit{Local Thermodynamic Equilibrium} (LTE) \citep{cristoforetti_local_2010}. Karena plasma LIBS bersifat sangat transien dengan densitas dan dimensi geometri spasial yang curam dari pusat (\textit{core}) hingga ke tepi luarnya, keadaan keseimbangan termodinamika secara utuh sangat mustahil untuk digapai. Namun dalam kondisi LTE, pada kompartemen elemen mikro (lokal) dan \textit{delay time} spesifik (sekitar skala $\mu$s), tumbukan inelastis antar elemen pengisi---terutama eksitasi lewat elektron---dipandang mengambil alih secara dominan atas transfer de-eksitasi energi radiatif spontan. Dalam kerangka asumsi ini, berbagai parameter fisis molekuler dan atomik cukup didekati melalui parameter termodinamika makroskopik seperti temperatur elektron ($T_e$) tunggal \citep{cristoforetti_local_2010}.

Keabsahan dominasi tumbukan tersebut dapat diverifikasi secara teoretis melalui pemenuhan kriteria McWhirter \citep{bultel_mcwhirter_2025, cristoforetti_local_2010}. Kriteria ini menetapkan prasyarat minimum untuk kerapatan elektron plasma ($n_e$) pada ambang temperatur tertentu:
\begin{equation} \label{eq:bab2_mcwhirter}
n_e \geq 1.6 \times 10^{12} \ T_e^{1/2} (\Delta E)^3 \, \text{cm}^{-3},
\end{equation}
dengan $\Delta E$ (dalam eV) merupakan selisih interval tingkat energi dari transisi optik atom dominan terlebar dari sistem yang diteliti.

% ======================================================================
% 2.3 Spektrum Emisi Plasma LIBS
% ======================================================================
\section{Spektrum Emisi Plasma LIBS}

\subsection{Distribusi Boltzmann}

Pada kondisi LTE, distribusi populasi tingkat energi eksitasi atom dalam sebuah tingkatan ionisasi dideskripsikan oleh Distribusi Boltzmann \citep{fujimoto_plasma_2004}:
\begin{equation} \label{eq:bab2_boltzmann}
n_i = \frac{n_z \, g_i \exp(-E_i / k_B T_e)}{Z_z(T_e)},
\end{equation}
di mana $n_i$ merupakan densitas populasi partikel pada level energi $i$, $n_z$ adalah total densitas partikel dari tingkatan ionisasi $z$, $g_i$ adalah degenerasi statistik tingkat energi, $E_i$ adalah energi tingkat eksitasi, $k_B$ adalah konstanta Boltzmann, dan $Z_z(T_e)$ adalah fungsi partisi (\textit{partition function}).

\subsection{Persamaan Saha}

Temperatur plasma ($T_e$) yang sama juga mengontrol keseimbangan populasi antar tingkatan ionisasi atom melalui Persamaan Saha \citep{fujimoto_plasma_2004}:
\begin{equation} \label{eq:bab2_saha}
\frac{n_e \, n_{z+1}}{n_z} = 2\,\frac{Z_{z+1}(T_e)}{Z_z(T_e)} \left(\frac{2\pi m_e k_B T_e}{h^2}\right)^{3/2} \exp\!\left(-\frac{E_{\mathrm{ion},z}}{k_B T_e}\right),
\end{equation}
dengan $n_{z+1}$ merupakan densitas partikel pada tingkat ionisasi satu tingkat lebih tinggi, $m_e$ adalah massa elektron, $h$ adalah konstanta Planck, dan $E_{\mathrm{ion},z}$ adalah energi ionisasi. Kedua persamaan ini menjadi fondasi utama \textit{forward model} untuk menghitung emisi kuantitatif pada kondisi plasma dengan parameter $T_e$ dan $n_e$ tertentu.

\subsection{Koefisien Einstein}

Populasi energi eksitasi yang kemudian meluruh ke tingkat energi yang lebih rendah akan meradiasikan foton emisi, yang intensitasnya ditakar oleh koefisien emisi spektral \citep{fujimoto_plasma_2004}:
\begin{equation} \label{eq:bab2_emiss}
j_\lambda = \frac{hc}{4\pi\lambda}\,A_{ki}\,n_k\,\phi_V(\lambda),
\end{equation}
di mana $A_{ki}$ adalah koefisien Einstein yang merepresentasikan probabilitas transisi spontan dari tingkat energi atas $k$ ke tingkat bawah $i$, $n_k$ adalah densitas populasi tingkat atas (dihitung via distribusi Saha--Boltzmann), dan $\phi_V(\lambda)$ adalah profil garis emisi.

\subsection{Pelebaran Garis Spektral}

Dalam pengukuran eksperimental, garis emisi yang secara teori bersifat monokromatik selalu mengalami pelebaran pada domain panjang gelombang. Dua mekanisme pelebaran yang dominan pada plasma LIBS adalah pelebaran Doppler dan pelebaran Stark.

Pelebaran Doppler disebabkan oleh pergerakan termal atom yang memancarkan foton, menghasilkan profil Gaussian dengan lebar \citep{griem_spectral_1974}:
\begin{equation} \label{eq:bab2_doppler}
\Delta\lambda_D = \lambda_{ki}\sqrt{\frac{8\ln 2\,k_B T_i}{m_k c^2}}.
\end{equation}

Pelebaran Stark disebabkan oleh interaksi medan listrik mikro dari partikel bermuatan di sekitar atom pemancar, menghasilkan profil Lorentzian dengan lebar \citep{griem_spectral_1974}:
\begin{equation} \label{eq:bab2_stark}
\Delta\lambda_S \approx 2\,\omega_{ki}(T_e)\,\frac{n_e}{n_{e,\mathrm{ref}}}.
\end{equation}

\subsection{Profil Voigt}

Dalam plasma LIBS, kedua mekanisme pelebaran di atas berlaku secara simultan. Superposisi dari profil Gaussian (Doppler) dan profil Lorentzian (Stark) menghasilkan profil gabungan yang dikenal sebagai Profil Voigt, yang dinyatakan melalui konvolusi keduanya dalam fungsi $\phi_V(\lambda)$ \citep{griem_spectral_1974}. Profil Voigt ini digunakan dalam term $\phi_V(\lambda)$ pada Persamaan~\ref{eq:bab2_emiss} untuk mensimulasikan bentuk garis emisi yang realistis.

\subsection{Spektrum Monoatomik dan Poliatomik}

Dari perspektif fisika dekomposisi spektral, profil emisi total poliatomik $S_{\text{poly}}(\lambda)$ yang dikoleksi via spektrometer merupakan superposisi linier dari emisi individual $S_{\text{mono}}^{(z)}(\lambda)$ untuk setiap elemen $z$. Spektrum monoatomik merangkum konstelasi puncak emisi individual suatu elemen spesifik dengan intensitas yang ditentukan oleh parameter plasma $T_e$ dan $n_e$.

Spektrum poliatomik campuran mewakili hasil jumlahan linier dari kontribusi setiap elemen yang dibobot oleh fraksi konsentrasi $c_z$:
\begin{equation} \label{eq:bab2_poly}
S_{\mathrm{poly}}(\lambda) = \sum_{z} c_z \, S_{\mathrm{mono}}^{(z)}(\lambda).
\end{equation}
Model dekomposisi eksplisit ini menjustifikasi konsep \textit{cross-attention} yang menghubungkan representasi spektrum monoatomik dengan observasi spektrum poliatomik dalam kerangka pembelajaran mendalam.

% ======================================================================
% 2.4 Kuantifikasi Unsur pada LIBS
% ======================================================================
\section{Kuantifikasi Unsur pada LIBS}

\subsection{Calibration Curve}
% TODO: Isi dengan penjelasan metode kalibrasi univariat/multivariat standar
%       pada LIBS, termasuk keterbatasannya pada matriks kompleks.

\subsection{Calibration-Free LIBS}
% TODO: Isi dengan penjelasan metode CF-LIBS, asumsi yang mendasarinya,
%       kelebihan dan keterbatasannya, serta hubungannya dengan asumsi LTE.

\subsection{X-ray Fluorescence (XRF) sebagai Data Referensi}
% TODO: Isi dengan penjelasan singkat prinsip XRF dan justifikasi penggunaannya
%       sebagai ground truth untuk validasi hasil kuantifikasi LIBS.

% ======================================================================
% 2.5 Deep Learning untuk Analisis Spektrum
% ======================================================================
\section{Deep Learning untuk Analisis Spektrum}

\subsection{Convolutional Neural Network (CNN)}

\textit{Convolutional Neural Network} (CNN) satu dimensi merupakan arsitektur \textit{deep learning} yang efektif untuk ekstraksi fitur lokal dari data sekuensial seperti spektrum LIBS. CNN bekerja melalui operasi konvolusi menggunakan jendela geser (\textit{kernel}) pada sumbu panjang gelombang untuk menangkap pola-pola lokal sinyal, seperti profil puncak emisi dan pelebaran garis Voigt \citep{he_deep_2016}. Aktivasi non-linear melalui fungsi \textit{Rectified Linear Unit} (ReLU) dan stabilisasi distribusi aktivasi melalui \textit{Batch Normalisation} \citep{ioffe_batch_2015} diterapkan pada setiap lapisan konvolusi untuk meningkatkan kapabilitas representasi model.

\subsection{Transformer}

Kompleksitas dependensi spektral pada rentang panjang gelombang yang luas (hingga lebih dari 35.000 kanal) membutuhkan mekanisme yang mampu menangkap relasi antar titik spektral secara simultan. Arsitektur \textit{Transformer} \citep{vaswani_attention_2017} menjawab kebutuhan ini melalui mekanisme \textit{Self-Attention} yang memetakan hubungan antar posisi spektral menggunakan matriks \textit{Query} ($\mathbf{Q}$), \textit{Key} ($\mathbf{K}$), dan \textit{Value} ($\mathbf{V}$):
\begin{equation} \label{eq:bab2_attention}
\mathrm{Attention}(\mathbf{Q},\mathbf{K},\mathbf{V}) = \mathrm{softmax}\!\left(\frac{\mathbf{Q}\mathbf{K}^\top}{\sqrt{d_k}}\right)\mathbf{V}.
\end{equation}
Skema \textit{Multi-Head Attention} memungkinkan model menangkap korelasi antar panjang gelombang dari berbagai subspace representasi secara paralel, sehingga mampu merekonstruksi konektivitas antar garis emisi dalam satu elemen, seperti relasi multiplet dari tingkat ionisasi yang sama.

\subsection{Encoder--Decoder}

Arsitektur \textit{Encoder--Decoder} memisahkan proses representasi masukan dan proses dekoding keluaran secara modular. Pada konteks penelitian ini, blok \textit{Encoder} memproses spektrum monoatomik $S_{\text{mono}}^{(z)}$ untuk menghasilkan representasi memori spesifik elemen, sementara blok \textit{Decoder} memproses spektrum poliatomik campuran.

\subsection{Cross-Attention}

Koneksi antara \textit{Encoder} dan \textit{Decoder} direkatkan melalui mekanisme \textit{Cross-Attention}. Pada mekanisme ini, blok \textit{Decoder} menghasilkan vektor $\mathbf{Q}_{\text{dec}}$ dari representasi spektrum poliatomik dan mencocokkannya dengan $\mathbf{K}_{\text{enc}}, \mathbf{V}_{\text{enc}}$ dari \textit{Encoder} yang menyimpan representasi spektrum monoatomik individual. Bobot \textit{cross-attention} yang dihasilkan merepresentasikan probabilitas kontribusi relatif setiap elemen monoatomik terhadap sinyal poliatomik campuran, sekaligus memungkinkan ekstraksi koefisien konsentrasi $c_z$ secara konsisten dengan prinsip dekomposisi spektral (Persamaan~\ref{eq:bab2_poly}).

\subsection{Domain Adaptation}

Model yang dilatih pada data sintetis (\textit{forward model} LTE) rentan terhadap penurunan kinerja saat diterapkan pada data eksperimental nyata akibat adanya \textit{domain shift}. Perbedaan ini meliputi efek instrumentasi, distorsi \textit{continuum background} residual, dan deviasi dari asumsi LTE ideal. Untuk mengatasi permasalahan ini, mekanisme adaptasi domain (\textit{domain adaptation}) melalui \textit{Gradient Reversal Layer} (GRL) \citep{ganin_domain_2016} diintegrasikan ke dalam arsitektur model. GRL bekerja dengan membalikkan arah gradien dari klasifier domain selama \textit{backpropagation}, sehingga mendorong \textit{encoder} untuk mengekstraksi fitur yang bersifat \textit{domain-invariant}---yaitu representasi spektral yang informatif untuk prediksi konsentrasi tanpa bias terhadap asal-usul domain data.

% ======================================================================
% 2.6 Evaluasi dan Validasi Data/Model
% ======================================================================
\section{Evaluasi dan Validasi Data/Model}

\subsection{Signal-to-Noise Ratio (SNR)}
% TODO: Isi dengan definisi SNR pada spektrum LIBS, cara perhitungannya,
%       dan kriteria minimum untuk kualitas data spektral yang dapat diterima.

\subsection{Peak Detection dan Kecocokan Garis Emisi}
% TODO: Isi dengan penjelasan metode identifikasi puncak emisi pada spektrum,
%       pencocokan terhadap basis data NIST ASD, dan kriteria kecocokan.

\subsection{Analisis Regresi Spektral}

Validasi spektrum sintetis merupakan tahap penting untuk memastikan bahwa spektrum hasil simulasi tidak hanya sesuai secara visual, tetapi juga konsisten secara kuantitatif dengan prinsip fisika plasma dan data referensi atomik. Dalam penelitian ini, validasi spektrum sintetis dilakukan melalui analisis regresi spektral yang mencakup validasi posisi garis emisi dan validasi distribusi intensitas garis.

Validasi posisi garis emisi dilakukan dengan membandingkan panjang gelombang hasil simulasi $\lambda_{\mathrm{sim}}$ terhadap panjang gelombang referensi dari \textit{NIST Atomic Spectra Database} (ASD), $\lambda_{\mathrm{NIST}}$. Hubungan keduanya dapat dianalisis melalui regresi linear:

\begin{equation}
\lambda_{\mathrm{sim}} = a\,\lambda_{\mathrm{NIST}} + b
\end{equation}

dengan nilai ideal $a \approx 1$ dan $b \approx 0$. Selain itu, galat posisi garis dapat dihitung menggunakan \textit{Root Mean Square Error} (RMSE):

\begin{equation}
\mathrm{RMSE}_{\lambda} =
\sqrt{\frac{1}{N}\sum_{i=1}^{N}
(\lambda_{\mathrm{sim},i}-\lambda_{\mathrm{NIST},i})^2}
\end{equation}

\subsection{Boltzmann Plot}

Validasi distribusi intensitas dilakukan menggunakan metode \textit{Boltzmann plot}, yang didasarkan pada hubungan:

\begin{equation}
\ln\left(
\frac{I\lambda}{g_k A_{ki}}
\right)
=
-\frac{E_k}{k_B T_e}+C
\end{equation}

di mana $I$ adalah intensitas garis emisi, $\lambda$ adalah panjang gelombang, $g_k$ adalah degenerasi tingkat atas, $A_{ki}$ adalah probabilitas transisi Einstein, $E_k$ adalah energi tingkat atas, $k_B$ adalah konstanta Boltzmann, dan $T_e$ adalah temperatur elektron plasma. Melalui regresi linear terhadap $E_k$, nilai kemiringan kurva dapat digunakan untuk memperkirakan temperatur plasma, sedangkan koefisien determinasi ($R^2$) menunjukkan tingkat konsistensi distribusi intensitas terhadap asumsi \textit{Local Thermodynamic Equilibrium} (LTE).

Selain itu, korelasi antara probabilitas transisi Einstein ($A_{ki}$) dan intensitas puncak hasil simulasi dapat dianalisis menggunakan korelasi Pearson atau regresi logaritmik untuk menilai apakah intensitas spektral meningkat secara konsisten terhadap probabilitas emisi atomik.

\subsection{RMSE, MAE, MAPE, dan Koefisien Determinasi ($R^2$)}
% TODO: Isi dengan definisi formal metrik evaluasi regresi (RMSE, MAE, MAPE, R^2),
%       relevansinya dalam konteks evaluasi prediksi konsentrasi unsur,
%       dan interpretasi nilai metrik dalam analisis kuantitatif LIBS.

\subsection{Analisis Attention Map}
% TODO: Isi dengan penjelasan tentang bagaimana bobot cross-attention dapat
%       divisualisasikan dan diinterpretasikan sebagai indikator kontribusi
%       setiap elemen monoatomik terhadap spektrum poliatomik,
%       serta relevansinya untuk validasi fisik model.

% ======================================================================
% 2.7 Penelitian Terkait dan Posisi Penelitian
% ======================================================================
\section{Penelitian Terkait dan Posisi Penelitian}

\subsection{Penelitian Berbasis Fisika Plasma}

Pengembangan arsitektur analitis untuk prediksi parameter komposisi unsur dari spektrum emisi plasma telah menyusuri berbagai lanskap metodologi. Arsitektur berbasis pemodelan fisika radiatif secara komprehensif dikembangkan dalam paket komputasi MERLIN \citep{favre_merlin_2025, favre_towards_2025}. Konsep pemodelan \textit{Radiative Transfer Equation} (RTE) model ini memungkinkan rekonstruksi spektrum yang konsisten secara termodinamika, meskipun kendala komputasi yang tinggi masih membatasi penerapannya untuk analisis \textit{real-time}.

\subsection{Penelitian Berbasis Machine Learning / Deep Learning}

Di sisi lain, pendekatan \textit{Machine Learning} (ML) dan \textit{Deep Learning} semakin banyak diterapkan untuk analisis kuantitatif LIBS. Riset terdahulu menunjukkan keunggulan pendekatan berbasis data dalam mengatasi pengaruh \textit{matrix effect} melalui model multivariat untuk kuantifikasi konsentrasi unsur \citep{hao_machine_2024, wang_simulation_2024}. Evolusi ke arah arsitektur hibrid CNN--Transformer secara berkesinambungan menyempurnakan kapabilitas analisis spektrum. Penelitian oleh \citet{liu_remote_2026} menunjukkan efektivitas CNN untuk menangkap pola puncak lokal yang dipadukan dengan blok \textit{Transformer} untuk menangkap relasi \textit{long-range dependencies} antar panjang gelombang pada spektrum poliatomik. Adaptasi arsitektur Informer berbasis \textit{ProbSparse Attention} pada data geologi juga divalidasi signifikan \citep{walidain_informer-based_2026, zhou_informer_2021}.

\subsection{Research Gap dan Kebaruan Penelitian}

Meskipun model \textit{Deep Learning} spektral tersebut mengungguli pendekatan konvensional, seluruh penelitian sebelumnya memperlakukan spektrum masukan sebagai sinyal komposit tunggal (\textit{single signal evaluation}). Belum ditemukan penelitian yang secara eksplisit memanfaatkan arsitektur dekomposisi struktural yang memisahkan representasi spektrum monoatomik individual dari spektrum poliatomik campuran. Mengintegrasikan relasi dekomposisi fisika ini melalui mekanisme \textit{Cross-Attention Encoder--Decoder} merupakan posisi kebaruan (\textit{novelty}) yang diusung dalam penelitian ini.